%!TEX program = xelatex
\documentclass[12pt,a4paper,openany,fleqn]{book}
\usepackage{fontspec, xunicode, xltxtra}
\usepackage{tikz}
\usepackage{subfig}
\usepackage{color}
\usepackage{graphicx} 
\usepackage{titlesec}
\usepackage{indentfirst} %首行缩进处理
\usepackage{geometry}  %使用geomtry包设置页面
\usepackage[colorlinks,linkcolor=black,anchorcolor=blue,citecolor=green]{hyperref} %使用超链接
\usepackage{wrapfig}
\usepackage[fleqn]{amsmath}
\usepackage{latexsym}
\usepackage{booktabs}


% 定义字体
\defaultfontfeatures{Mapping=tex-text}

% 定义了两种用颜色%%%%%
%\definecolor{foot}{RGB}{39,129,214}  %定义页脚的文字颜色
\newfontfamily\kai{KaiTi}
	
% 定义一个新的命令来强调文字下划线红色
\newcommand{\emp}[1]{\color{red}{#1}}
\newcommand{\tabincell}[2]{\begin{tabular}{@{}#1@{}}#2\end{tabular}}

% 使用fancyhdr定义页面的设置
\usepackage{fancyhdr} 
\pagestyle{fancy}
\lhead{}
\chead{}
\rhead{}


\setlength{\parskip}{12pt} %%设置段落间的距离
\setlength{\parindent}{2em}  %%首行缩进两个字符

%% 用fncychap设置章节的样式
%\ChNameVar{\LARGE\bfseries}
%\ChTitleVar{\LARGE\bfseries}{第\,\thechapter\,章}

%% 用[titlesec]设置章节的格式
\titleformat{\part}[display]{\Huge\bfseries}{第\,\thepart\,篇}{1em}{}
\titleformat{\chapter}{\LARGE\bfseries}{第\,\thechapter\,章}{1em}{}
\titleformat{\section}{\Large\bfseries}{第\,\thesection\,节}{1em}{}


% 设置文本字体和断句方式
\setmainfont{MicrosoftYaHei}  
\XeTeXlinebreaklocale "zh"
\XeTeXlinebreakskip = 0pt plus 1pt minus 0.1pt
\linespread{1.25}  %%行距扩大为默认的1.28倍

% 重新命名目录
\renewcommand\contentsname{\textbf{目录}}
\renewcommand\figurename{图}
\renewcommand\tablename{表}
\renewcommand{\thefootnote}{\alph{footnote}}


%文档开始
\begin{document}

% 目录页面
\thispagestyle{empty}
\begin{titlepage}
\begin{flushright} 

% 设置文档的名称、写文档的作者内容

%% 画线包括书名
\vspace*{3cm} %%起始位置添加3cm的空行

\begin{tikzpicture}
\draw[color=cyan,line width=3pt] 
({0.4\textwidth},0) --(\textwidth,0);
 	
\end{tikzpicture} 
\\[0.8cm]

{\color{cyan} {\Huge {汇添富企业版手册}}} \\[0.8cm]

% 画一条双线
\begin{tikzpicture}
\draw[color=cyan,line width=3pt] 
({0.4\textwidth},0) --(\textwidth,0);
 	
\end{tikzpicture} \\[1cm]

{\Large 版本号：1.0} \\[5cm]
 
{\large 汇添富企业版产品组}\\[1cm]

{\large 2017年04月10日} \\[1cm]
%Powered By \LaTeX


\end{flushright}
\end{titlepage}


%%扉页的内容
\clearpage
\thispagestyle{empty}

本文档仅用于产品经理、商务拓展经理和其它相关人员学习和查询汇添富企业版所有产品的功能和相应的操作，{\emp{请各位妥善保存，如非必要尽量不要随意分发}}。如果您在使用过程中发现本文描述有误或者没有及时更新，您可以把问题发送邮件至 \href{mailto:zhengyouxin@htffund.com}{企业版产品组邮箱}，我们会及时的更新文档。


%% 感谢名单
\clearpage 

\thispagestyle{empty}
\begin{center}
{\Large 感谢}\\[1cm]

 企业版WEB开发人员: \\
{\kai {王志刚、董奇、谭作军、付志辉、彭春荣}} \\[1cm]

 企业版APP开发人员: \\
{\kai {胡宗权、廖伟伟、林师权、廖倩文}} \\[1cm]

 测试组人员: \\
{\kai {李彧奇、方炼炼}} \\[1cm]

 企业版产品组人员: \\
{\kai {梁晨、周娟、钟凌戈、张俊玲、郑又鑫}} \\[1cm]

同时感谢各位领导和其他组同事支持和帮助

\end{center}

% 前言部分
\clearpage
\begin{center}  %% 前言页面
{\Large 前言}
\\[.6cm]

\end{center}


早在2014年，汇添富基金就推出了企业版理财的客户端，提供了可以在Windows平台安装的企业现金管理工具，企业客户可以通过客户端进行现金宝的申购、赎回（普通取现和快速取现）、收益查看，同时还提供高端理财产品的申购赎回等功能。\par

2015年7月底，汇添富上线企业版\url{https:qy.99fund.com}，开始将客户端的功能迁移到WEB平台，提供更完整、更方便的企业现金管理工具和平台，并将原有的客户端客户逐渐迁移到新的平台。汇添富企业版通过官方网站提供更丰富的功能和产品，并持续地进行用户体验的优化，时至今日，汇添富企业版经过数十次的迭代优化，已经能很好的满足企业客户的需求。\par

2015年12月，汇添富发布企业版APP，提供了资产查询和产品查看的功能，同时微信服务号正式上线，用户可以绑定企业版账户，通过微信服务号查询资产和产品。我们后期再不断的迭代优化企业版APP，目前经过数次版本发布，企业版APP已经提供完整的功能，包括用户管理、交易、业务管理等。企业用户也可以通过APP完成绝大多数的日常业务，您可以点击\href{https://qy.99fund.com/appScan/appOnline.htm}{这里}
查看更多的信息,除此之外，我们也提供微信服务号和H5站点等产品。\par

为了更好的帮助产品经理、商务拓展经理、客服和客户了解企业版和企业版的产品，产品组整理了当前企业版产品的功能，特别的编写了本文档，方便用户更好的了解企业版产品。\par

本文档分为三大篇，分别是产品篇、用户篇和系统管理篇，特别是产品篇和用户篇重点介绍了企业版各产品的功能和操作内容，并且提供了相关问题的定义和解释，对产品经理、商务拓展经理和用户有很大帮助，系统管理篇主要介绍了企业版产品某些关系到用户的个性设置和系统管理功能，方便客服人员进行查阅。\par

希望本文档能够有效帮助到使用企业版的用户，祝企业版越来越好！\par

\begin{flushright}
	企业版产品组\\
	2016年11月23日
\end{flushright}


% 目录部分
\clearpage
\thispagestyle{empty}
\tableofcontents


% 第一篇产品
\part{产品篇}
\label{par:product}

% 这里是对产品篇的简介

% 第一章：企业版web
\clearpage
\chapter{企业版介绍} % (fold)
\label{cha:inst_intro}
\noindent{
	本章节包含以下内容：\\[14pt]
	$\Rightarrow$ 企业版产品简介 \\[12pt]
	$\Rightarrow$ 企业版WEB产品（参见章节\ref{sec:inst_web}）\\[12pt]
	$\Rightarrow$ 企业版APP产品（参见章节\ref{sec:inst_app}）\\[12pt]
	$\Rightarrow$ 企业版微信服务号（参见章节\ref{sec:inst_wx}）\\[12pt]
	$\Rightarrow$ 企业版H5产品（参见章节\ref{sec:inst_h5}）\\[12pt]
}

企业版是汇添富推出的面向企业客户的投资理财平台。向企业客户提供货币基金、专户理财和基金产品，企业客户通过互联网线上的方式在企业版注册、认购、申购、赎回、查看资产等功能。\par

企业版产品包括了WEB、APP、H5、微信服务号，其中WEB和APP是企业版核心产品，H5和微信服务号是辅助产品。WEB和APP提供完整的功能，包括新用户注册、登录、查看产品、申认购产品、查看资产等功能。H5和微信服务号提供部分功能，诸如登录、查看资产等功能。\par

企业版不同的产品具备了不同的功能集合，在不同的场景下向客户提供及时需要的解决方案，因此我们需要对不同的产品有不同的定位，面向不同的用户群体，所以在功能设计上有不同的侧重点。\par

具体的产品定位、目标用户、功能设计，可以参考如下的列表。\par

%  表格1：产品功能表
\begin{table}[htbp]
\caption{产品定位及其目标用户}
\begin{tabular*}{\textwidth}{lccc}
	\toprule  %--------------------------------------------------------------------------------------
	企业版产品      & 产品定位                  & 目标用户                    & 产品功能           \\
	\midrule  %--------------------------------------------------------------------------------------
    企业版WEB       & 核心产品   & \tabincell{c}{业务人员\\财务总监\\老板}   & 交易、查看、产品、内容 \\
    %\hline
    企业版APP       & 核心产品   & \tabincell{c}{财务总监\\老板\\业务人员}   & 交易、查看、产品、内容 \\
    %\hline
    企业版H5        & 辅助产品   & \tabincell{c}{业务人员\\财务经理}         & 查看、产品、推荐、活动 \\
    %\hline
    企业版微信服务号 & 辅助产品   & \tabincell{c}{业务人员}                  & 资产查看、推送、推荐    \\
	\bottomrule %------------------------------------------------------------------------------------
\end{tabular*}
\end{table}

关于业务人员\footnote{指具体业务操作人员包括经办员、复核员和高级用户}、财务总监\footnote{具体的投资的决策者}、老板\footnote{公司负责人对投资有决定权}，参见对其的备注描述。企业版需要通过产品的展示、推荐内容、活动内容向投资决策者传递投资产品的信息，主要的产品渠道是APP、微信服务号以及企业版WEB。同时，向具体的业务人员提供产品渠道用于购买赎回产品。\par

% 第二章：企业版产品
\clearpage
\chapter{企业版产品} 
\label{cha:inst_product}
% 2.1节企业版web
\section{企业版WEB} 
\label{sec:inst_web}
企业版WEB是汇添富企业版的核心和拳头产品，也是汇添富最早向客户推出的企业投资平台之一。\href{https://qy.99fund.com}{汇添富企业版}是业内最早发布的企业理财平台。企业客户可以在企业版购买货币基金、基金产品和高端理财，因此企业版主要分为现金宝、高端理财、基金以及其它功能。 \par

% 插入企业版首页图片
\begin{figure}[htbp!]
  \centering
  \fbox{
  \includegraphics[width=0.9\textwidth]{picture/homepage.png}}
  \caption{汇添富企业版首页}
\end{figure}

汇添富企业版主要面向具体业务人员以及财务总监，业务人员可以在汇添富企业版购买和赎回产品，也可以查看产品的信息。财务总监主要是可以作为管理员对用户进行管理，也可以查看具体的产品信息以决定是否购买。企业版WEB是汇添富向企业客户提供的重点产品，提供了企业版最完备的功能，包括注册新用户、完善资料、登录、查看资产、购买产品、赎回产品等等。


% 2.2节企业版app
\section{企业版APP} 
\label{sec:inst_app}
% 插入APP下载二维码
\begin{wrapfigure}{L}{5.6cm}
	\includegraphics[width=0.27\textwidth]{picture/app_home.png}
	\label{fig:app_home}
	%\caption{app页面}
\end{wrapfigure}

企业版APP是汇添富向企业客户推出的移动端产品，用户可以通过企业版APP完成交易等操作。企业版APP也是汇添富企业版的核心产品，用户可以扫描\hyperref[fig:app_home]{\emp{左方的二维码}}下载。\par

企业版APP于2015年12月正式上线，2016年2月份正式登陆AppStore。用户可以在AppStore和应用宝安卓市场\href{http://sj.qq.com/myapp/detail.htm?apkName=com.fund.instmobile}{企业现金宝}下载程序。\par

左下方是APP的图示，目前企业版APP的功能包括：注册新用户、完善资料、用户角色管理、密码管理、购买和赎回基金、高端理财产品的购买赎回等。另外，提供复核业务推送功能，经办员发起业务后推送至复核员，复核员可以在APP完成业务复核。\par

企业版APP是企业版布局在移动互联网的产品，为客户在移动端提供解决方案，也包括了除企业代发之外的几乎所有功能。并且易于获得和使用，建议客户和我司的商务拓展经理下载使用。并可以将使用中遇到的问题反馈给\href{mailto:zhengyouxin@htffund.com}{企业版产品组邮箱}。我们会持续的优化和迭代产品。


% 2.3节企业版微信服务号
\clearpage
\section{企业版微信服务号} 
\label{sec:inst_wx}
% 微信服务号介绍页面
\begin{wrapfigure}{L}{5.6cm}
	\includegraphics[scale=0.7]{picture/scanWechat.jpg}
	\label{fig:scanWechat}
\end{wrapfigure}
微信服务号是对企业版产品的补充，可以让用户在不方便访问企业版WEB和在没有安装企业版APP时可以查询其持有的资产和交易记录（显示最近10条），您可以扫描\hyperref[fig:scanWechat]{\emp{左边二维码}}关注服务号，登录关联企业账号后就可以查看关联账户的资产和交易记录。\par

微信服务号提供了一些基本的功能，主要包括账户绑定微信、切换绑定账户、查看当前绑定的账户信息、查看企业信息、查看企业持有资产、查看交易记录、查看现金宝产品、查看高收益产品、以及可以通过微信服务号申请预约码。\par

\begin{wrapfigure}{R}{0pt}
	\includegraphics[width=0.45\textwidth]{picture/weixin.pdf}
	\label{fig:weixin}
\end{wrapfigure}

基于微信的方便性和易用性，随时可以获取企业的重要数据，给用户带来了便利性。另外，微信服务号基于微信平台，具有跨平台的功能。\par

出于安全性考虑，以及企业版的交易功能比较复杂，且具有非常规的流程，所以交易功能不予提供。\par

后续为了更方便的服务用户，经办员发起业务的消息会通过微信服务号推送至复核员，并且提供H5页面以方便复核员查看并复核业务。\par

另外，微信服务号可以通过H5页面的方式访问并查看现金宝、高端理财和基金的产品。



\clearpage
% 2.4节企业版微信服务号
\section{企业版H5} 
\label{sec:inst_h5}
企业版H5是汇添富推出的手机浏览器访问的网站，有基本的功能，主要是注册新用户、完善资料、查看持有资产等。客户可以通过访问\href{mqy.99fund.com}{企业版H5网站}查看资产和注册为新用户。H5网站的主要功能是引流吸引新用户，商务拓展经理在介绍企业版时，如果不方便用电脑网络访问企业版，可以使用手机浏览H5网站，向客户展示企业版的特征。 \par

\vspace{1.5cm}

\noindent {\Large 总结} \par

企业版的产品形态多样，并且不同的产品提供了不同的功能，适用于不同的使用场景。建议商务拓展经理在向客户推荐企业版时，可以推荐所有产品并重点展示企业版WEB和企业版APP，这两个产品具有了常用的完备的功能，并且是企业版的核心产品。\par

具体的不同产品的定位和功能可以参考表格\ref{tab:function}。





% 第二篇用户管理
\clearpage
\part{用户篇}
\label{par:users}

% 这里是对用户篇的简介

% 第三章：用户注册
\clearpage
\chapter{用户注册} 
\label{cha:reg}
% 本章介绍企业版客户的注册

% 3.1小节 新用户在企业版web注册
\section{企业版WEB注册} 
\label{sec:reg_web}
企业版WEB是企业版最主要的产品，提供最完整的功能，建议用户使用企业版WEB进行注册为企业版用户。用户在企业版注册会直接生成为企业的交易账号和基金账号，并且创建管理员角色。\par

\noindent{
	注册前需要准备的资料：}\\
	\begin{enumerate}
		\item 管理员手机号码\footnote{建议财务总监或者财务经理}（可以接受短信）
		\item 企业的证件（营业执照、组织机构代码证、税务登记证）
		\item 企业的证件（统一社会信用代码）
		\item 法人的身份证件
		\item 管理员的身份证件信息
		\item 银行开户证明文件
	\end{enumerate}

企业用户注册前请先准备好以上的证件和信息，包括证件的复印件或者扫描件（后续过程需要上传），用户可以访问\href{https://qy.99fund.com/instReg/register.htm}{这里注册}。\par

% 第一步：填写手机号码
\noindent{第1步：填写手机号码并设置登录密码}\\
\begin{figure}[htbp!]
  \centering
  \fbox{
  \includegraphics[width=0.9\textwidth]{picture/reg1.png}}
  \caption{填写注册信息}
\end{figure}

% 第二步：填写企业信息
\noindent{第2步：填写企业信息}\\
\begin{figure}[htbp!]
  \centering
  \fbox{
  \includegraphics[width=0.9\textwidth,height=6.5cm]{picture/reg2.png}}
  \caption{填写企业信息}
\end{figure}

% 第三步：填写法人信息
\noindent{第3步：填写法人代表信息}\\
\begin{figure}[htbp!]
  \centering
  \fbox{
  \includegraphics[width=0.9\textwidth]{picture/reg3.png}}
  \caption{填写法人代表信息}
\end{figure}

% 第四步：填写管理员信息
\noindent{第4步：填写管理员信息}\\
\begin{figure}[htbp!]
  \centering
  \fbox{
  \includegraphics[width=0.9\textwidth]{picture/reg4.png}}
  \caption{填写管理员信息}
\end{figure}

填写完以上信息后，确认无误点击下一步。完成信息的提交，注册的第一阶段完成。这时候在汇添富企业版已经生成企业账户，并根据第1步的手机号码生成管理员账号。\par

% 第五步：注册完成
\noindent{第5步：完成注册}\\
\begin{figure}[htbp!]
  \centering
  \fbox{
  \includegraphics[width=0.9\textwidth]{picture/regs.png}}
  \caption{完成注册}
\end{figure}

\noindent{\large 需要注意的事项:}\\
\begin{enumerate}
	\item 管理员登录密码6-18位
	\item 证件类型为可选：营业执照\&统一社会信用代码
	\item 选择营业执照需要填写其它两个证件号码
	\item 允许一个手机号码注册多个企业账户
	\item 不允许同一证件号码\footnote{指营业执照或者统一社会信用代码，组织机构代码和税务登记证号码同一不影响}注册
	\item 同一证件号码要注册多账户的情况可以通过加后缀实现
  \item 以上第2步至第5步均在一个页面完成
\end{enumerate}
\par
企业版完成注册后，可以使用手机号码（管理员角色）登录企业版，并进一步完成完善资料,参见章节\ref{sec:update_web}。\par


% 3.2小节 新用户在企业版APP注册
\section{企业版APP注册} 
\label{sec:reg_app}
本节介绍在企业版APP注册为新用户。 \par

% 3.3小节 新用户在H5注册
\section{企业版H5注册} 
\label{sec:reg_h5}
本节介绍在企业版H5网站上注册为新用户。 \par



% 第四章：用户迁移
\clearpage
\chapter{用户迁移} 
\label{cha:transfer}
% 第四章-本章介绍企业版客户的登录

% 4.1小节 客户端用户迁移到企业版
\section{客户端用户迁移} 
\label{sec:transfer_client}
本节介绍客户端客户迁移至企业版。 \par

% 4.2小节 柜台用户迁移到企业版
\section{柜台客户迁移} 
\label{sec:transfer_otc}
本节介绍柜台客户迁移到企业版。 \par

% 第五章：用户登录
\clearpage
\chapter{用户登录} 
\label{cha:login}
% 本章介绍企业版客户的登录

% 5.1小节 用户在企业版web登录
\section{企业版WEB登录} 
\label{sec:login_web}
点击\href{https://qy.99fund.com/instLogin/login.htm}{这里}进入企业版登录页面。登录页面初始的信息为手机号码、登录角色、登录密码.\par

% 插入初始的登录页面
\begin{figure}[htbp!]
	\centering
	\fbox{
	\includegraphics[width=0.9\textwidth]{picture/web_login1.png}
	}
	\caption{汇添富登录页面}
\end{figure}

用户输入手机号码后，系统会自动根据手机号码判断注册的企业，显示营业执照号码，并显示该手机号码注册的用户角色。用户根据营业执照号码和角色选定后，输入登录密码后点击登录。\par

\begin{figure}[htbp!]
	\centering
	\fbox{
	\includegraphics[width=0.9\textwidth]{picture/web_login2.png}
	}
	\caption{汇添富登录页面}
\end{figure}

用户登录完成后，会直接进入到我的账户页面。如果手机号码被多个企业注册了，在营业执照号会显示所有该号码注册的营业执照号码，选择了该营业执照号码后会显示在该企业注册的用户角色。

% 5.2小节 用户在企业版APP登录
\section{企业版APP登录} 
\label{sec:login_app}

% 插入APP登录图片
\begin{figure}[htbp!]
	\centering
	\fbox{
	\includegraphics[width=0.80\textwidth]{picture/app_login.pdf}
	}
	\caption{汇添富登录页面}
\end{figure}

APP登录方式和WEB一致，在输出手机号码后会根据号码检查其注册情况，列出其注册的营业执照号码，选择其需要登录的执照号码后，会显示对应企业账户的用户角色，输入登录密码后登录。\par

打开APP后默认为登录页面，其初始登录信息和WEB一致，登录信息包括手机号码、登录角色和登录密码。\par


% 5.3小节 用户在H5登录
\section{企业版H5登录} 
\label{sec:login_h5}
% h5登录页面
\begin{wrapfigure}{L}{0.5\textwidth}
	\includegraphics[width=0.45\textwidth]{picture/h5_login.pdf}
	\label{fig:h5_login}
	\caption{H5登录页面}
\end{wrapfigure}

H5登录页面和APP不一样，初始的信息包含了手机号码、营业执照号码、登录角色和登录密码输入框，输入手机号码后会显示该手机注册的企业的营业执照号码，选择营业执照号码后选择对应的用户角色，输入登录密码后完成登录。\par


% 5.4小节 用户在微信服务号登录
\section{企业版微信服务号登录} 
\label{sec:login_wx}
用户可以扫描\ref{fig:scanWechat}微信服务号二维码关于企业版微信服务号，或者如以下图示搜索并关注企业版微信公众服务号。 \par

\begin{figure}[htbp!]
  \centering
  \fbox{
  \includegraphics[width=0.9\textwidth]{picture/wx_gz.pdf}}
  \caption{关注企业版微信服务号}
\end{figure}

% 5.5小节 用户登录时忘记密码
\section{忘记登录密码} 
\label{sec:login_passwd}
本节介绍忘记登录密码后的处理方式。 \par


% 第六章：完善资料
\clearpage
\chapter{完善资料} 
\label{cha:update}
% 本章介绍企业版客户的注册

% 6.1小节 新用户在企业版web完善资料
\section{企业版WEB完善资料} 
\label{sec:update_web}
本节介绍在企业版web完善资料。 \par

% 6.2小节 新用户在企业版APP完善资料
\section{企业版APP完善资料} 
\label{sec:update_app}
本节介绍在企业版APP完善资料。 \par

% 6.3小节 新用户在H5完善资料
\section{企业版H5完善资料} 
\label{sec:update_h5}
本节介绍在企业版H5网站完善资料。 \par


% 第七章：用户角色
\clearpage
\chapter{用户角色} 
\label{cha:role}
% 本章介绍企业版客户的用户角色内容

% 7.1小节 用户角色说明
\section{用户角色说明} 
\label{sec:role_intro}
本节介绍企业版的用户角色。 \par

% 7.2小节 管理用户角色
\section{用户角色管理} 
\label{sec:role_manage}
用户角色管理是指管理员对其它用户角色的创建、修改和删除等操作。\par


% 7.2.1小节 创建用户角色
\subsection{创建用户角色} 
\label{sub:role_add}
管理员可以创建非管理员用户角色。 \par

如无任何非管理员用户，在【我的账户】页面的有创建经办员、复核员和高级用户的提醒。\par

\begin{figure}[htbp!]
  \centering
  \fbox{
  \includegraphics[width=0.9\textwidth]{picture/role_tips.png}}
  \caption{创建用户角色提醒}
\end{figure}

点击【创建角色】跳转至角色管理页面创建用户，也可以点击【角色说明】查看不同角色的说明。\par

点击【我的账户】-【用户管理】进入用户角色管理页面，点击【新增角色】创建用户角色。\par

\begin{figure}[htbp!]
  \centering
  \fbox{
  \includegraphics[width=0.9\textwidth]{picture/role_add.png}}
  \caption{新增用户角色}
\end{figure}

\begin{figure}[htbp!]
  \centering
  \fbox{
  \includegraphics[width=0.9\textwidth]{picture/role_create.png}}
  \caption{新增用户角色}
\end{figure}

% 7.2.2小节 修改用户角色
\subsection{修改用户角色} 
\label{sub:role_modi}
本节介绍如何修改用户角色。 \par

% 7.2.3小节 删除用户角色
\subsection{删除用户角色} 
\label{sub:role_del}
本节介绍删除用户角色。 \par


% 第八章：密码管理
\clearpage
\chapter{密码管理} 
\label{cha:passwd}
% 本章介绍企业版客户的用户密码管理

% 8.1小节 修改用户密码
\section{修改用户密码} 
\label{sec:passwd}
本节介绍如何修改用户的密码，\par


% 8.2小节 重置非管理员用户角色密码
\section{重置非管理员用户密码} 
\label{sec:passwd_reset}

% 8.2.1小节 重置登录密码
\subsection{重置登录密码} 
\label{sub:passwd_login_reset}
本节介绍管理员重置非管理员用户的登录密码重置。\par

% 8.2.2小节 重置交易密码
\subsection{重置交易密码} 
\label{sub:passwd_trade_reset}
本节介绍管理员重置非管理员的交易密码。 \par



% 第三篇交易
\clearpage
\part{交易篇}
\label{par:trade}

% 这里是对交易篇的简介

% 第九章：现金宝交易
\clearpage
\chapter{现金宝交易} 
\label{cha:trade_mac}
% 本章介绍企业版现金宝交易

% 9.1小节 现金宝充值
\section{现金宝充值} 
\label{sec:trade_mac_recharge}
本节介绍如何申购现金宝。 \par

% 9.2小节 现金宝赎回
\section{现金宝赎回} 
\label{sec:trade_redemp}

% 9.2.1小节 快速取现
\subsection{快速取现} 
\label{sub:trade_mac_quick_redemp}
本节介绍现金宝的快速取现。 \par

% 9.2.2小节 普通赎回
\subsection{普通赎回} 
\label{sub:trade_mac_normal_redemp}
本节介绍现金宝普通取现。 \par

% 9.3小节 现金宝收益查询
\section{现金宝收益查询} 
\label{sec:trade_mac_record}
本节介绍现金宝交易记录查询和对账单功能. \par



% 第十章：高收益交易
\clearpage
\chapter{高收益交易} 
\label{cha:trade_vip}
% 本章介绍企业版高端理财交易

% 10.1小节 高端理财认购
\section{高端理财认购} 
\label{sec:trade_vip_buy}
本节介绍高端理财产品的认购。 \par

% 10.2小节 高端理财申购
\section{高端理财申购} 
\label{sec:trade_vip_charge}

% 10.2.1小节 预约购买
\subsection{预约购买} 
\label{sub:trade_vip_prebuy}
本节介绍高端理财产品的预约购买。 \par

% 10.2.2小节 立即购买
\subsection{立即购买} 
\label{sub:trade_vip_qbuy}
本节介绍直接购买高端理财。 \par

% 10.2.3小节 申请预约吗
\subsection{申请预约吗} 
\label{sub:trade_vip_precode}
本节介绍如何申请高端理财的预约码。 \par

% 10.3小节 高端理财赎回
\section{高端理财赎回} 
\label{sec:trade_vip_redemp}

% 10.3.1小节 高端理财预约赎回
\subsection{高端理财预约赎回} 
\label{sub:trade_vip_redemp_pre}
本节介绍高端理财的预约赎回。 \par

% 10.3.2小节 高端理财赎回设置
\subsection{高端理财赎回设置} 
\label{sub:trade_vip_redemp_setting}
本节介绍高端理财预约赎回设置。 \par


% 第十一章：基金交易
\clearpage
\chapter{基金交易} 
\label{cha:trade_fund}
% 本章介绍企业版基金交易

% 11.1小节 基金认购
\section{基金认购} 
\label{sec:trade_fund_rebuy}
本节介绍基金的认购。 \par

% 11.2小节 基金申购
\section{基金申购} 
\label{sec:trade_fund_buy}
本节介绍基金的申购。 \par

% 11.3小节 基金赎回
\section{基金赎回} 
\label{sec:trade_fund_redemp}
本节介绍基金的赎回。 \par


% 第十二章：交易管理
\clearpage
\chapter{交易管理} 
\label{cha:trade_manage}
% 本章介绍企业版交易管理

% 12.1小节 交易复核
\section{交易复核} 
\label{sec:trade_manage_check}
本节介绍交易复核。 \par

% 12.2小节 交易取消
\section{交易取消} 
\label{sec:trade_manage_drop}
本节介绍取消交易。 \par

% 12.3小节 交易撤单
\section{交易撤单} 
\label{sec:trade_manage_cancel}
本节介绍撤回交易. \par


% 第十三章：交易记录和对账单
\clearpage
\chapter{交易记录和对账单} 
\label{cha:trade_record}
% 本章介绍企业版交易记录和对账单

% 13.1小节 资产查看
\section{资产查看} 
\label{sec:trade_record_cap}
本节介绍如何查看资产。 \par

% 13.2小节 交易记录查看
\section{交易记录查看} 
\label{sec:trade_record_check}
本节介绍如何查看交易记录。 \par

% 13.3小节 对账单
\section{对账单} 
\label{sec:trade_record_list}
本节介绍如何查看月度和年度对账单。 \par


% 第四篇系统管理
\clearpage
\part{系统管理篇}
\label{par:admin}

% 这里是对系统管理篇的简介


% 第五篇参考内容
\clearpage
\part{参考篇}
\label{par:ref}

% 这里是对参考篇的简介
\clearpage
\chapter{参考表格} % (fold)
\label{cha:ref_table}

\section{产品功能表} % (fold)
\label{sec:function_table}
%  表格1：产品功能表
\begin{table}[htbp]
\caption{产品功能表}
\label{tab:function}
\begin{tabular*}{\textwidth}{@{\extracolsep{\fill}}lcccc}
	\toprule  %--------------------------------------------------------------------------------------
	功能                   &WEB            &APP                 &H5               &微信服务号      \\
	\midrule  %--------------------------------------------------------------------------------------
	注册                  &$\surd$         &$\surd$            &$\surd$           &$\times$       \\
	完善资料               &$\surd$         &$\surd$            &$\surd$           &$\times$       \\
	登录                  &$\surd$         &$\surd$            &$\surd$           &$\surd$        \\
	查看资产               &$\surd$         &$\surd$            &$\surd$           &$\surd$        \\
	查看明细资产            &$\surd$         &$\surd$            &$\surd$           &$\surd$        \\
	发起交易               &$\surd$         &$\surd$            &$\times$          &$\times$        \\
	复核业务               &$\surd$         &$\surd$            &$\times$          &$\times$        \\
	查看交易记录            &$\surd$         &$\surd$            &$\times$          &$\surd$        \\
	收益查询               &$\surd$         &$\surd$            &$\times$           &$\surd$        \\
	查看对账单             &$\surd$         &$\times$           &$\times$           &$\times$        \\
	用户管理               &$\surd$         &$\surd$            &$\surd$            &$\times$       \\
	消息推送               &$\times$        &$\surd$            &$\times$           &$\times$       \\
	\bottomrule
\end{tabular*}
\end{table}





% 文档结束
\end{document}